% !PRINTED 254
\weberpage{254}

\begin{center}
\S.\,66.\\
\gesperrt{Simplicity of the Group $E$.}
\end{center}

The first question in a more thorough investigation of the group $E$ is that of a possibly existing normal divisor. It can be proved that, apart from two quite simple exceptions, no such normal divisor exists.

Let us assume that $G$ is a normal divisor of $E$ which contains at least one substitution different from the identity,
\weq{1}{A = \begin{pmatrix} \alpha, & \beta \\ \gamma, & \delta \end{pmatrix}}
By the definition of a normal divisor, if $T$ is an arbitrary element of $E$, then
\weq{2}{TAT^{-1}}
is also contained in $G$, and it is to be shown that in this way, and by composition of such substitutions, one can derive all the elements of $E$, from which it then follows that $G$ must be identical with $E$. For this, however, it suffices to establish the occurrence of the generating substitutions $A_0$, $B_h$ [\S.\,65, (8)] in $G$.

We proceed to this step by step.

1) In $G$ there occurs a substitution of the form
\weq{3}{\begin{pmatrix} 0, & \beta \\ \gamma, & \delta \end{pmatrix}}
To show this, we form according to (2):
\weq{4}{\begin{pmatrix} \xi, & 1 \\ -1, & 0 \end{pmatrix}\begin{pmatrix} \alpha, & \beta \\ \gamma, & \delta \end{pmatrix}\begin{pmatrix} 0, & -1 \\ 1, & \xi \end{pmatrix} = \begin{pmatrix} \beta\xi+\delta, & \beta\xi^2+(\delta-\alpha)\xi-\gamma \\ -\beta, & -\beta\xi+\alpha \end{pmatrix}.}
Now if $\beta$ is different from zero, one can determine $\xi$ from $\beta\xi+\delta=0$ and the goal is attained. But if $\beta=0$, then $\gamma=0$, $\delta-\alpha=0$ do not hold simultaneously, since otherwise $A$ would be the identical substitution; one can therefore take $\xi$ in such a way that $(\delta-\alpha)\xi-\gamma$ is different from zero. Then by (4) one has formed a substitution $A$ in $G$ in which $\beta$ is not zero, and to which the formula (4) can now be applied once more, in order to make $\alpha$ zero.

2) In $G$ there occurs a substitution of the form
\[
\begin{pmatrix} 0, & 1 \\ -1, & \xi \end{pmatrix}
\]
%CONT>>
